% ════════════════════════════════════════════════════════════
% SYNTHÈSE DES RÉSULTATS
% ════════════════════════════════════════════════════════════
\clearpage
\sectionspeciale{Synthèse des résultats}

Cette section récapitule l'ensemble des résultats du dimensionnement, séance par séance. Elle permet une lecture rapide des dimensions, sections d'aciers et vérifications retenues sur chaque élément structurel du bâtiment.

\subsection*{Séance 1 --- Prédimensionnement}

\begin{tabularx}{\linewidth}{Xc}
\toprule
\textbf{Élément} & \textbf{Dimensions retenues} \\
\midrule
Dalle nervurée (infrastructure)              & $h = 12$~cm \\
Plancher-dalle (superstructure)              & $h = 29$~cm (par itération) \\
Poutre secondaire (infrastructure)           & $20 \times 55$~cm$^{2}$ \\
Poutre principale (infrastructure)           & $30 \times 70$~cm$^{2}$ \\
Poteau de superstructure                     & $50 \times 50$~cm$^{2}$ \\
Poteau d'infrastructure                      & $60 \times 60$~cm$^{2}$ \\
\bottomrule
\end{tabularx}

\subsection*{Séance 2 --- Descente de charges (poteau B2 au 2\textsuperscript{e} sous-sol)}

\begin{tabularx}{\linewidth}{Xc}
\toprule
\textbf{Grandeur} & \textbf{Valeur} \\
\midrule
Surface tributaire (avec continuité $\times 1{,}15$) & $S^* = 44{,}85$~m$^{2}$ \\
Charges permanentes cumulées $\Sigma G$              & $2\,460$~kN \\
Charges d'exploitation cumulées $\Sigma Q$ (réduites) & $1\,238$~kN \\
Effort normal ELU au pied du poteau $N_{Ed}$         & $\approx 5\,179$~kN \\
\bottomrule
\end{tabularx}

\subsection*{Séance 3 --- Semelles de fondations}

\begin{tabularx}{\linewidth}{Xc}
\toprule
\textbf{Solution} & \textbf{Dimensions / armatures} \\
\midrule
\textbf{Solution sur 2 pieux} (retenue)        & Pieux $\varnothing$ 0{,}70~m, entraxe 1{,}75~m \\
~~~~~Semelle de liaison                       & $2{,}55 \times 0{,}80 \times 1{,}10$~m \\
~~~~~Armatures inférieures $A_{inf}$           & $11$ HA 25 ($\approx 54$~cm$^{2}$) \\
~~~~~Armatures supérieures $A_{sup}$           & $A_{inf}/10 \approx 5{,}4$~cm$^{2}$ \\
~~~~~Cadres verticaux et horizontaux           & HA 14 e=10 cm \\
\textbf{Solution alternative semelle isolée} & $2{,}70 \times 2{,}70 \times 0{,}65$~m \\
~~~~~Armatures dans chaque sens                & $11$ HA 25 \\
\bottomrule
\end{tabularx}

\subsection*{Séance 4 --- Dalle sur 4 côtés (1\textsuperscript{er} sous-sol)}

\begin{tabularx}{\linewidth}{Xc}
\toprule
\textbf{Grandeur} & \textbf{Valeur} \\
\midrule
Dimensions de la dalle                          & $3{,}05 \times 5{,}85$~m, épaisseur 12 cm \\
Coefficient $\alpha = l_x/l_y$                  & $0{,}521$ (travaille dans deux sens) \\
Moment max en travée sens $x$ ($M_{tx}$)        & $\approx 8{,}6$~kN.m/ml \\
Moment max en travée sens $y$ ($M_{ty}$)        & $\approx 6{,}4$~kN.m/ml \\
Sections d'acier en travée                       & $A_{tx} = 3{,}29$~cm$^{2}$/ml ; $A_{ty} = 2{,}25$~cm$^{2}$/ml \\
\bottomrule
\end{tabularx}

\subsection*{Séances 5 \& 6 --- Poutre continue (file B)}

\begin{tabularx}{\linewidth}{Xc}
\toprule
\textbf{Grandeur} & \textbf{Valeur} \\
\midrule
Géométrie de la poutre                              & $30 \times 70$~cm sur 2 travées (6{,}10 et 5{,}90 m) \\
Moment maximal sur appui $M_{B2}$                   & $453$~kN.m \\
Moment maximal en travée                            & ELU $\approx 280$ à $300$~kN.m \\
Armatures sur appui B2 (méthode Caquot)             & $6$ HA 25 \\
Armatures en travée                                  & $3$ HA 25 (avec arrêt de barres) \\
Armatures sur appui B1 (rive)                       & $3$ HA 25 \\
Cadres transversaux                                  & HA 8 espacement variable \\
\bottomrule
\end{tabularx}

\subsection*{Séance 7 --- Escalier}

\begin{tabularx}{\linewidth}{Xc}
\toprule
\textbf{Grandeur} & \textbf{Valeur} \\
\midrule
Géométrie (règle de Blondel)                & Giron $g = 30$~cm, contremarche $h = 17{,}5$~cm \\
Inclinaison de la paillasse                  & $\alpha = 30^\circ$ \\
Épaisseur de la paillasse (projetée)         & $22{,}6$~cm \\
Charge ELU                                    & $P_u = 13{,}2$~kN/m$^{2}$ \\
Armatures principales paillasse               & $A_{inf} = 3{,}85$~cm$^{2}$/ml \\
Armatures de répartition                      & HA 8 e=20~cm \\
Béquet --- armatures                          & $A_s = 1{,}57$~cm$^{2}$/ml \\
\bottomrule
\end{tabularx}

\subsection*{Séance 8 --- Mur pignon (épaisseur 20 cm)}

\begin{tabularx}{\linewidth}{Xc}
\toprule
\textbf{Vérification} & \textbf{Résultat} \\
\midrule
Effort normal ELU au RDC                         & $N_{Ed} = 634$~kN/ml \\
Capacité résistante ELU $N_{Rd}$                  & $678$~kN/ml \\
Vérification de stabilité                         & $N_{Ed} < N_{Rd}$ \quad $\checkmark$ \\
Contrainte max ELU $\sigma_{max,3}$               & $3{,}17$~MPa \\
\midrule
\textbf{Console courte au R+1}                    & \\
~~Tirant principal $A_{tirant}$                   & $4{,}78$~cm$^{2}$ \\
~~Aciers répartis horizontaux $A_r$               & $1{,}20$~cm$^{2}$ \\
~~Aciers répartis verticaux $A_v$                 & $9{,}28$~cm$^{2}$ \\
~~Anti-éclatement $A_t$                           & $3{,}01$~cm$^{2}$ \\
\bottomrule
\end{tabularx}

\subsection*{Séance 9 --- Mur de sous-sol (file C)}

\begin{tabularx}{\linewidth}{Xc}
\toprule
\textbf{Élément} & \textbf{Valeur} \\
\midrule
\textbf{Soutènement (méthode forfaitaire FDP 18-717)} & \\
~~Coefficient de poussée                              & $K_0 = 0{,}50$ \\
~~Poussée ELU max (à $-5{,}40$ m)                    & $73{,}1$~kN/m$^{2}$ \\
~~Moment isostatique travée basse $M_{0,1}$          & $51{,}7$~kN.m/ml \\
~~Moment en travée $M_{t,1}$                          & $38{,}8$~kN.m/ml \\
~~Moment en travée $M_{t,2}$                          & $14{,}03$~kN.m/ml \\
~~Armatures travée 1 / appui / travée 2               & $5{,}40$ / $4{,}34$ / $1{,}90$ cm$^{2}$/ml \\
\midrule
\textbf{Poutre cloison (méthode bielles-tirants)}     & \\
~~Force tirant $F_t$                                  & $447$~kN \\
~~Force bielle $F_b$ (avec 2 bielles)                 & $1\,112$~kN ($\sigma_b = 11{,}1 \leq 12{,}75$~MPa) \\
~~Section tirant $A_t$                                & $10{,}3$~cm$^{2}$ \\
\bottomrule
\end{tabularx}

\subsection*{Séance 10 --- Calcul au feu (R120)}

\begin{tabularx}{\linewidth}{Xc}
\toprule
\textbf{Vérification} & \textbf{Résultat} \\
\midrule
\textbf{Poutre principale 30$\times$70}                & \\
~~Moment résistant en travée $M_{Rt}$                  & $85{,}55$~kN.m \\
~~Moment résistant sur appui B2 $M_{Ra}(B_2)$           & $699{,}52$~kN.m \\
~~Moment résistant sur appui B1 $M_{Ra}(B_1)$           & $392{,}55$~kN.m \\
~~Vérification $M_{Rt} + (M_{Re}+M_{Rw})/2 \geq M_0$    & $631{,}59 \geq 413$~kN.m \quad $\checkmark$ \\
\midrule
\textbf{Dalle nervurée 12 cm}                          & \\
~~Capacité portante sens $x$                           & $\bar{p}_x = 17{,}31$~kN/ml \\
~~Capacité portante sens $y$                           & $\bar{p}_y = 16{,}65$~kN/ml \\
~~Charge sollicitante                                   & $p = g + 0{,}8\,q = 11$~kN/ml \\
~~Vérification                                          & $\bar{p}_x, \bar{p}_y \geq p$ \quad $\checkmark$ \\
\bottomrule
\end{tabularx}

\subsection*{Séances 11 \& 12 --- Plancher-dalle}

\begin{tabularx}{\linewidth}{Xc}
\toprule
\textbf{Grandeur} & \textbf{Valeur} \\
\midrule
\textbf{Moments fléchissants}                          & \\
~~Moment isostatique de référence $M_0$                & $459$~kN.m \\
~~Moment sur appui A (rive)                            & $115$~kN.m \\
~~Moment sur appui B (intermédiaire)                   & $275{,}4$~kN.m \\
~~Moment en travée AB                                  & $365$~kN.m \\
~~Moment en travée BC                                  & $345{,}8$~kN.m \\
~~Section maximale d'acier (travée AB, bandes pot.)    & $\approx 11$~cm$^{2}$ \\
\midrule
\textbf{Poinçonnement au poteau B2}                    & \\
~~Cisaillement sollicitant $v_{Ed}$                    & $0{,}67$~MPa \\
~~Cisaillement résistant $v_{Rd,c}$                     & $0{,}46$~MPa \\
~~Conclusion                                            & Armatures de poinçonnement requises \\
~~Section dimensionnante $A_{sw}/S_r$                   & $34{,}48$~cm$^{2}$ (à 0{,}56 m) \\
~~Choix d'armatures                                     & $61$ HA $10$ ou $42$ HA $12$ \\
~~Vérification au nu du poteau                          & $v_{Ed} = 1{,}69 \leq v_{Rd,max} = 3{,}6$~MPa \quad $\checkmark$ \\
\bottomrule
\end{tabularx}

\subsection*{Bilan des dimensions retenues pour l'ouvrage}

\begin{tabularx}{\linewidth}{Xc}
\toprule
\textbf{Élément structurel} & \textbf{Dimensions / Sections} \\
\midrule
Dalle nervurée (infrastructure)             & $h = 12$~cm \\
Plancher-dalle (superstructure)             & $h = 29$~cm \\
Poutre secondaire                           & $20 \times 55$~cm$^{2}$ \\
Poutre principale                           & $30 \times 70$~cm$^{2}$ \\
Poteau de superstructure                    & $50 \times 50$~cm$^{2}$ \\
Poteau d'infrastructure                     & $60 \times 60$~cm$^{2}$ \\
Mur pignon                                  & épaisseur 20 cm \\
Mur de sous-sol (file C)                    & épaisseur 20 cm \\
Pieux de fondation                          & $\varnothing$ 0{,}70~m \\
Semelle sur 2 pieux                         & $2{,}55 \times 0{,}80 \times 1{,}10$~m \\
Semelle superficielle (alternative)         & $2{,}70 \times 2{,}70 \times 0{,}65$~m \\
\bottomrule
\end{tabularx}

\vspace{1cm}

\subsection*{Conformité réglementaire}

L'ensemble des éléments structurels du bâtiment a été dimensionné selon les Eurocodes (EN 1990 à EN 1997) et le FDP 18-717. Les vérifications suivantes ont été menées et satisfaites :

\begin{itemize}
  \item Stabilité aux \textbf{ELU} (état limite ultime) sous combinaisons fondamentales pour tous les éléments porteurs ;
  \item Aptitude au service \textbf{ELS} (état limite de service) pour les fondations sur pieux ;
  \item Stabilité au \textbf{feu R120} (2~heures) pour la poutre principale et la dalle nervurée ;
  \item \textbf{Poinçonnement} du plancher-dalle au droit du poteau intermédiaire ;
  \item \textbf{Dispositions constructives} (sections minimales, espacements, enrobages) sur l'ensemble des éléments en béton armé ;
  \item Action du \textbf{vent} sur le mur pignon et sa console au R+1 ;
  \item \textbf{Poussée des terres} sur le mur de sous-sol et répartition des charges des poteaux par poutre cloison.
\end{itemize}

L'ouvrage répond aux exigences de \textbf{stabilité} et de \textbf{sécurité} fixées par les normes en vigueur, et constitue une base solide pour la phase de conception détaillée.

